\chapter{RK4 time integration}
\label{chap:rk4-time-integration}

\begin{quote}
Spec dir:
\passthrough{\lstinline!specs/2026-05-02-phase-6-rk4-integrator/!}.
Version: \passthrough{\lstinline!0.6.0!}.
\end{quote}

\subsection{What this gives you}\label{what-this-gives-you}

A higher-order time integrator (classic 4-stage Runge--Kutta) and a
\textbf{generic dispatcher} that lets you pick which integrator to use:

\begin{lstlisting}[language={C++}]
solve::integrate(psi, rhs, dt, n_steps, IntegratorTag{});
\end{lstlisting}

with \passthrough{\lstinline!IntegratorTag!} being either
\passthrough{\lstinline!solve::ExplicitEuler!} (Phase 5) or
\passthrough{\lstinline!solve::RK4!} (new in Phase 6).

\passthrough{\lstinline!solve::diffuse!} is now templated on the
integrator tag so you can swap between Euler and RK4 with one extra
argument:

\begin{lstlisting}[language={C++}]
solve::diffuse(psi, D, dt, n_steps);                    // Euler (default)
solve::diffuse(psi, D, dt, n_steps, solve::RK4{});      // RK4
\end{lstlisting}

RK4 is \textbf{more accurate} at the same \passthrough{\lstinline!dt!}
(4 orders better --- the time-error per step is
\passthrough{\lstinline!O(dt\^5)!} vs Euler's
\passthrough{\lstinline!O(dt\^2)!}) and has a \textbf{larger stability
bound} for the heat equation
(\passthrough{\lstinline!D dt sum(1/h\_a\^2) ≤ 0.6963!} vs Euler's
\passthrough{\lstinline!0.5!}). The cost is 4 RHS evaluations per step
instead of 1.

\subsection{API}\label{api}

Headers:
\href{../../../include/gridcalc/solve/Integrator.hpp}{\passthrough{\lstinline!include/gridcalc/solve/Integrator.hpp!}}
(one-stop aggregator),
\href{../../../include/gridcalc/solve/ExplicitEuler.hpp}{\passthrough{\lstinline!include/gridcalc/solve/ExplicitEuler.hpp!}},
\href{../../../include/gridcalc/solve/RK4.hpp}{\passthrough{\lstinline!include/gridcalc/solve/RK4.hpp!}},
\href{../../../include/gridcalc/solve/Diffusion.hpp}{\passthrough{\lstinline!include/gridcalc/solve/Diffusion.hpp!}}.

\begin{lstlisting}[language={C++}]
namespace gridcalc::solve {

// Integrator tag types -- empty structs that carry compile-time properties.
struct ExplicitEuler { static constexpr double diffusionCFLLimit = 0.5;    };
struct RK4           { static constexpr double diffusionCFLLimit = 0.6963; };

// Generic dispatcher (overloaded on tag).
template <class Rhs>
void integrate(core::Field<double>& psi, Rhs&& rhs, double dt, int n_steps, ExplicitEuler);

template <class Rhs>
void integrate(core::Field<double>& psi, Rhs&& rhs, double dt, int n_steps, RK4);

// Phase-5 thin wrapper -- equivalent to integrate(..., ExplicitEuler{}).
template <class Rhs>
void explicitEuler(core::Field<double>& psi, Rhs&& rhs, double dt, int n_steps);

// Diffusion driver -- now templated on integrator tag.
template <class Integrator = ExplicitEuler>
void diffuse(core::Field<double>& psi, double D, double dt, int n_steps,
             Integrator tag = {});

}  // namespace gridcalc::solve
\end{lstlisting}

\subsubsection{Picking an integrator}\label{picking-an-integrator}

{\def\LTcaptype{none} % do not increment counter
\begin{longtable}[]{@{}
  >{\raggedright\arraybackslash}p{(\linewidth - 2\tabcolsep) * \real{0.8088}}
  >{\raggedright\arraybackslash}p{(\linewidth - 2\tabcolsep) * \real{0.1912}}@{}}
\toprule\noalign{}
\begin{minipage}[b]{\linewidth}\raggedright
Want
\end{minipage} & \begin{minipage}[b]{\linewidth}\raggedright
Use
\end{minipage} \\
\midrule\noalign{}
\endhead
\bottomrule\noalign{}
\endlastfoot
Simplest, smallest per-step cost &
\passthrough{\lstinline!ExplicitEuler\{\}!} (Phase 5 default) \\
Best accuracy per step & \passthrough{\lstinline!RK4\{\}!} \\
Larger time step (closer to RK4's \passthrough{\lstinline!\~1.39×!}
looser CFL) & \passthrough{\lstinline!RK4\{\}!} \\
Pure linear / Hamiltonian / oscillatory dynamics &
\passthrough{\lstinline!RK4\{\}!} (lower phase error) \\
One-pass throwaway integration of a smooth IC & Either; default Euler is
fine \\
\end{longtable}
}

Both integrators take the same RHS callable signature:
\passthrough{\lstinline!Field<double> rhs(const Field<double>\&)!}.
SFINAE-validates; mismatch trips a
\passthrough{\lstinline!static\_assert!}.

\subsubsection{CFL stability check}\label{cfl-stability-check}

\passthrough{\lstinline!solve::diffuse!} validates the von-Neumann bound
\passthrough{\lstinline!D · dt · sum\_a(1/h\_a\^2) ≤ Tag::diffusionCFLLimit!}
before integration begins and throws
\passthrough{\lstinline!std::invalid\_argument!} if violated. The error
message names the integrator type, the offending values, and the
applicable limit.

\passthrough{\lstinline!solve::integrate!} and
\passthrough{\lstinline!solve::explicitEuler!} do \textbf{not} check CFL
--- they're generic over the RHS, which means they cannot know the
relevant stability bound. If you call them directly with a custom RHS,
you own the stability argument.

\subsection{Worked example: heat equation with
RK4}\label{worked-example-heat-equation-with-rk4}

\begin{lstlisting}[language={C++}]
#include <cmath>
#include <gridcalc/core/Field.hpp>
#include <gridcalc/core/Grid.hpp>
#include <gridcalc/solve/Diffusion.hpp>
#include <gridcalc/solve/Integrator.hpp>

using gridcalc::core::Field;
using gridcalc::core::Grid;
using gridcalc::core::Vec3d;
using gridcalc::solve::diffuse;
using gridcalc::solve::RK4;

constexpr double kPi = 3.14159265358979323846;

const int N = 32;
const double L = 2.0 * kPi;
const double h = L / N;
const double D = 0.1;
const double dt = 0.05;
const int n_steps = 100;

Grid grid(N, N, N, Vec3d(h, h, h));

Field<double> psi(grid, [](double x, double y, double z) {
  return std::sin(x) * std::sin(y) * std::sin(z);
});

// One-line switch from Euler to RK4: pass the integrator tag.
diffuse(psi, D, dt, n_steps, RK4{});

// psi now matches the analytical exp(-3 D T) sin(x) sin(y) sin(z) to
// within ~0.5% relative max-norm error -- 10x tighter than Phase 5's
// Euler at the same parameters.
\end{lstlisting}

\subsection{Worked example: custom RHS (no CFL
gating)}\label{worked-example-custom-rhs-no-cfl-gating}

\begin{lstlisting}[language={C++}]
#include <gridcalc/diff/Laplacian.hpp>
#include <gridcalc/solve/Integrator.hpp>

// Reaction-diffusion: dpsi/dt = D laplacian(psi) - k psi
auto rhs = [D, k](const Field<double>& f) {
  Field<double> out = gridcalc::diff::laplacian(f);
  const std::size_t N_ = out.getSize();
  double* o = out.data();
  const double* p = f.data();
  for (std::size_t n = 0; n < N_; ++n) {
    o[n] = D * o[n] - k * p[n];
  }
  return out;
};

// RK4 makes this much more accurate than Euler at the same dt.
gridcalc::solve::integrate(psi, rhs, dt, n_steps, gridcalc::solve::RK4{});
\end{lstlisting}

\subsection{\texorpdfstring{What this does \emph{not} do
(yet)}{What this does not do (yet)}}\label{what-this-does-not-do-yet}

\begin{itemize}
\tightlist
\item
  \textbf{No adaptive time stepping.} Embedded RK methods (RK4(5),
  Dormand--Prince) arrive only if a downstream phase needs them.
\item
  \textbf{No symplectic / structure-preserving integrators.} Heat
  equation is dissipative; symplectic schemes are for Hamiltonian flows.
\item
  \textbf{No multistep methods (Adams--Bashforth, BDF).} Implicit BDF
  arrives in Phase 19.
\item
  \textbf{No persistent scratch pool for RK4.} Each step allocates
  \textasciitilde6 fields; Phase 20 swaps in a reusable scratch buffer.
\item
  \textbf{Boundary conditions: still periodic-only.} Phase 18's
  Dirichlet/Neumann policies will need their own integrator-test
  coverage; tracked in \passthrough{\lstinline!STATUS.md!} ``Open /
  deferred items.''
\end{itemize}

For the why-RK4-is-order-4 derivation and the stability-bound
calculation, see
\href{../../developer-note/notes/phase-6-rk4-integrator.md}{\passthrough{\lstinline!docs/developer-note/notes/phase-6-rk4-integrator.md!}}.
